\documentclass[10pt]{article}

% set smaller margins
\usepackage[a4paper,margin=1in]{geometry}

% use standard math symbols
\usepackage{amssymb}

% for \text inside formulas and for \eqref
\usepackage{amsmath}

% enable full XeLaTeX power
\usepackage{xltxtra}

% select widely available fonts
\setmainfont{Palatino}

% set paragraph margins
\setlength{\parindent}{0pt}
\setlength{\parskip}{6pt plus 2pt minus 1pt}

% prevent overfull lines
\setlength{\emergencystretch}{3em}

% for \includegraphics
\usepackage{graphicx}

% for exact placement of figures
\usepackage{float}

% for hyperlinks
\usepackage[colorlinks,urlcolor=blue,filecolor=blue,linkcolor=black,citecolor=black]{hyperref}
\usepackage[all]{hypcap}
\urlstyle{same}

\begin{document}

\section{Minimal Texts.js Sample}

This document shows what syntax must be supported in texts.js to be compatible
with Markdown (\href{http://johnmacfarlane.net/pandoc/}{pandoc} dialect).

\section{Headings}

\subsection{Subsection Heading}

\subsubsection{Subsubsection}

\section{Text Formatting}

You can use \emph{emphasis} and \textbf{strong} emphasis. Inline \texttt{code}
will use monspaced font.

\begin{verbatim}
Code blocks use monospaced font as well and preserve line
breaks
\end{verbatim}

\section{Math}

Math formulas can be used inside paragraph like $E=mc^2$ or on a separate line like the following one.

\begin{equation}
1+\frac{q^2}{(1-q)}+\frac{q^6}{(1-q)(1-q^2)}+\cdots =
\prod_{j=0}^{\infty}\frac{1}{(1-q^{5j+2})(1-q^{5j+3})},
\quad\quad \text{for $|q|<1$}
\end{equation}

\section{Footnotes}

Footnotes are placed anywhere\footnote{Yes right here.} in the source file but
appear at the bottom of the page.

\section{Hyperlinks}

You can add arbitrary hyperlinks. Link \href{http://www.texts.io/}{text} can
differ from it's URL or be the same as the URL: \url{http://www.google.com/}.

You can include e-mail links:
\href{mailto:sheremetyev@gmail.com}{sheremetyev@gmail.com}. Local files can be
\href{run:basic.pdf}{referenced} as well.

\section{Images}

\begin{figure}[H]
\centering
\includegraphics{Texts_Logo.png}
\end{figure}

\section{Tables}

\begin{table}[H]
\centering
\begin{tabular}{ | l | l | l | p{8cm} |}
\hline
Day & Min Temp & Max Temp & Summary \\ \hline
Monday & 11C & 22C & A clear day with lots of sunshine.  However, the strong breeze will bring down the temperatures. \\ \hline
Tuesday & 9C & 19C & Cloudy with rain, across many northern regions. Clear spells across most of Scotland and Northern Ireland, but rain reaching the far northwest. \\ \hline
Wednesday & 10C & 21C & Rain will still linger for the morning. Conditions will improve by early afternoon and continue throughout the evening. \\
\hline
\end{tabular}
\end{table}

\section{Lists}

\begin{itemize}

\item First bulleted item.

\item Second bulleted item.

    \begin{itemize}

    \item Subitem.

        \begin{itemize}

        \item Subsubitem.

        \end{itemize}

    \end{itemize}

\end{itemize}


\begin{enumerate}

\item First numbered item.

\item Second numbered item.

    \begin{enumerate}

    \item Subitem.

        \begin{enumerate}

        \item Subsubitem.

        \end{enumerate}

    \end{enumerate}

\end{enumerate}

\section{Quoting}

\begin{quotation}
Quoted text can include only inline text styles.
\end{quotation}

\begin{quotation}
Headings shouldn’t be used.
\end{quotation}

\section{Symbol Escaping}

No-break space (Unicode U+00A0) will be converted to tilde symbol for TeX.

\XeTeX{} doesn't need \TeX{} ligatures for quotation marks, em-dash and other
typographic symbols—they can be ‘used’ “directly”. Euro symbol can be used for
as less as €10. Ellipsis… in the text should be a Unicode character.

Should \TeX{} ligatures be ``used''—they will be printed <<as is>>.

\TeX{} special symbols are automatically escaped: \%, \$, \{, \}, \_, \#,
\&, \textbackslash, \textasciitilde, \textasciicircum.

HTML entites don't have to be escaped: \&, <, >.

\section{Future Extensions}

\begin{itemize}

\item Definition lists.

\item Titles for images, tables and formulas.

\item Cross-references.

\item Bibliographic citations.

\item Comments.

\item Custom labels for blocks.

\end{itemize}

\end{document}
